\section{Sissejuhatus}
\label{sissejuhatus}

\subsection{Probleemi taust ja motivatsioon}

Metsade süsinikuvaru ja selle muutuste hindamine on oluline nii kliimapoliitika kujundamisel kui ka ökosüsteemide seisundi seires. Lamapuit (surnud puit maapinnal) on üks kriitilisemaid metsade süsinikuringe, struktuurse mitmekesisuse ja elurikkuse komponente. Teaduskirjanduses peetakse lamapuitu üheks olulisemaks metsade biodiversiteedi indikaatoriks, kuna see pakub elupaika ja toitu tuhandetele saproksüülsetele liikidele, olles seega ökosüsteemi tervise otsene peegeldaja \parencite{lassauce_deadwood_2011, stokland_biodiversity_2012}. Vaatamata sellele on lamapuidu ruumiline ja kvantitatiivne hindamine jätkuvalt keerukas, tuginedes Eestis ja mujal suures osas piiratud ulatusega välitöödele ning statistilistele üldistustele, mis ei paku piisavat ruumilist resolutsiooni \parencite{marchi_airborne_2018}.

Aerolaserskaneerimise (Light Detection and Ranging – LiDAR) tehnoloogia on viimastel kümnenditel revolutsioneerinud metsanduslikku kaugseiret, võimaldades hinnata puistu struktuuri kolmemõõtmeliselt. Lisaks ökoloogilisele väärtusele on lamapuidu täpne kaardistamine kriitilise tähtsusega praktilises metsahalduses. Näiteks on lamapuidu maht oluline sisendparameeter metsatulekahjude leviku modelleerimisel, kuna maapinnal asuv kuiv surnud puit moodustab olulise osa kütusekoormusest \parencite{shin_morphological_2024, dietenberger_accurate_2025}. Operatiivne lamapuidu tuvastamine on asendamatu tööriist ka tormikahjude ulatuse kiireks hindamiseks, võimaldades planeerida koristustöid ja hinnata metsade kahjustusastet vahetult pärast loodusõnnetusi \parencite{jarron_detection_2021}.

Eestis kogutakse aerolaserskaneerimise andmeid riikliku seire raames (ALS IV ring), kuid üleriigilise lennu keskmine punktitihedus (umbes 2,1 p/m\textsuperscript{2}) seab piirid väiksemate objektide, nagu üksikute tüvede, detailsele tuvastamisele \parencite{maa-ja_ruumiamet_als_nodate}. Samas on Maa- ja Ruumiamet alustanud madallendude teostamist, kus punktitihedus ulatub üle 18 p/m\textsuperscript{2}, mis loob eeldused masinnägemise meetodite rakendamiseks lamapuidu automaatseks segmenteerimiseks. Senine praktika on näidanud, et suurim takistus selliste mudelite arendamisel on kvaliteetsete märgistatud treeningandmete puudus.

\subsection{Töö eesmärk ja uurimisküsimused}

Käesoleva magistritöö peamine eesmärk on demonstreerida kaasaegsete masinõppe meetodite (klassifitseerimine ja segmenteerimine) teostatavust lamapuidu tuvastamiseks Eestis, kasutades selleks kõrge resolutsiooniga aerolaserskaneerimise andmeid. Erinevalt varasematest statistilistest lähenemistest keskendub käesolev töö objektipõhisele analüüsile tihedas punktipilves. Töö oluliseks väljundiks on ka esimese standardiseeritud lamapuidu treeningandmestiku loomine Eesti metsatüüpide baasil, mis võimaldab tulevikus arendada mudeleid üleriigiliseks rakendamiseks.

Töös püütakse vastata järgmistele uurimisküsimustele:
\begin{itemize}
    \item Millise täpsusega on võimalik klassifitseerida ja segmenteerida lamapuitu Maa- ja Ruumiameti madallendude (18+ p/m\textsuperscript{2}) punktipilve andmetest?
    \item Kas sügavõppe-põhised masinnägemise mudelid on võimelised eristama lamapuitu muust maapinnalähedasest taimestikust Eesti komplekssetes metsaoludes?
    \item Milline on loodud treeningandmestiku potentsiaal edasisteks arendusteks, et liikuda üleriigilise automaatse seiresüsteemi suunas?
\end{itemize}

\subsection{Uurimisülesanded}

Töö eesmärgi saavutamiseks on kavandatud järgmised uurimisülesanded:
\begin{itemize}
    \item analüüsida teaduskirjanduse põhjal parimaid praktikaid lamapuidu tuvastamiseks tihedates punktipilvedes;
    \item teostada kõrge punktitihedusega LiDAR-andmete eeltöötlus ja normaliseerimine;
    \item märgistada ja ette valmistada esinduslik treeningandmestik lamapuidu objektidest;
    \item treenida ja valideerida masinõppe mudeleid lamapuidu klassifitseerimiseks ja segmenteerimiseks valitud testaladel;
    \item hinnata mudelite täpsust ning analüüsida peamisi vigu (nt segiminek rattaroobastega või madala põõsarindega).
\end{itemize}